\begin{table}[htbp]
\centering
\small
\caption{Spatial Boundary Discontinuity Estimates Across Locational Deliverability Area (LDA) Borders}
\label{tab:table4_rdd}
\adjustbox{max width=\textwidth}{%
\begin{tabular}{lcccc}
\toprule
 & (1) & (2) & (3) & (4) \\
Dependent Variable: & Residential Rate & Industrial Rate & $\ln(\text{Mfg Emp})$ & Q1 Energy Burden \\
 & (cents/kWh) & (cents/kWh) & (NAICS 325/331) & (\% Income) \\
\midrule
Inside Constrained LDA & 3.214*** & 2.845*** & -0.184*** & 1.984*** \\
 & (0.428) & (0.395) & (0.042) & (0.245) \\
Distance to LDA Border (km) & 0.0073 & -0.0040 & -0.0018 & 0.0045 \\
 & (0.0066) & (0.0066) & (0.0014) & (0.0041) \\
\midrule
Bandwidth & $\pm 25$ km & $\pm 25$ km & $\pm 25$ km & $\pm 25$ km \\
Quarter Fixed Effects & Yes & Yes & Yes & Yes \\
Observations & 1,440 & 1,440 & 1,440 & 1,440 \\
$R^2$ & 0.815 & 0.810 & 0.725 & 0.815 \\
\bottomrule
\end{tabular}%
}
\begin{minipage}{\linewidth}
\vspace{1ex}
\footnotesize \textit{Notes:} Standard errors clustered at the county level in parentheses. * $p<0.10$, ** $p<0.05$, *** $p<0.01$. Sample restricted to contiguous county borders separating transmission-constrained LDAs (BGE, DOM) from unconstrained zones within a 25 km geographic bandwidth. Column (1) documents a discontinuous +3.214 cents/kWh retail tariff jump ($t = 7.51$). Column (3) estimates a statistically significant -18.4\% manufacturing employment contraction ($\beta = -0.184, t = -4.38$), corresponding to -420 manufacturing jobs per border county.
\end{minipage}
\end{table}
